\section{Parte 3: Redes Virtuales y Seguridad Perimetral}

\subsection{3.1 Diseño y Segmentación de Red Virtual (\texttt{VNet-01})}
La arquitectura de red define el perímetro de aislamiento para los servicios de cómputo y datos. Se aprovisionó la red virtual corporativa \textbf{\texttt{VNet-01}} con un espacio de direccionamiento privado \texttt{10.0.0.0/16}, segmentada en dos subredes especializadas:
\begin{itemize}[leftmargin=1.5em]
  \item \textbf{\texttt{Subnet-Frontend} (\texttt{10.0.1.0/24})}: Destinada a puntos de entrada, balanceadores y servicios públicos.
  \item \textbf{\texttt{Subnet-Backend} (\texttt{10.0.2.0/24})}: Destinada a microservicios internos, bases de datos y orquestación sin exposición a Internet.
\end{itemize}

\begin{azureterminal}[az network vnet create \& az network vnet subnet create]
$ az network vnet create \
    --resource-group "rg-bootcamp-part03-networking" \
    --name "VNet-01" \
    --address-prefix "10.0.0.0/16" \
    --subnet-name "Subnet-Frontend" \
    --subnet-prefix "10.0.1.0/24" \
    --location "eastus"

$ az network vnet subnet create \
    --resource-group "rg-bootcamp-part03-networking" \
    --vnet-name "VNet-01" \
    --name "Subnet-Backend" \
    --address-prefix "10.0.2.0/24"

{
  "addressPrefix": "10.0.2.0/24",
  "id": ".../virtualNetworks/VNet-01/subnets/Subnet-Backend",
  "name": "Subnet-Backend",
  "provisioningState": "Succeeded"
}
\end{azureterminal}

\subsection{3.2 Filtrado Granular con Grupos de Seguridad de Red (NSG)}
Para mitigar accesos no autorizados y controlar el tráfico perimetral, se configuró el Network Security Group \textbf{\texttt{nsg-bootcamp-sec}} asociado a las interfaces de red. Se establecieron reglas determinísticas de control de flujo con prioridades estrictas:

\begin{table}[h!]
\centering
\small
\begin{tabularx}{\textwidth}{lccccX}
\toprule
\textbf{Regla} & \textbf{Prioridad} & \textbf{Acceso} & \textbf{Puerto} & \textbf{Protocolo} & \textbf{Propósito Arquitectónico} \\
\midrule
\texttt{Allow-HTTP} & 100 & Allow & 80 & TCP & Tráfico Web no seguro (redirección) \\
\texttt{Allow-HTTPS} & 110 & Allow & 443 & TCP & Tráfico Web Seguro cifrado TLS 1.3 \\
\texttt{Allow-SSH} & 120 & Allow & 22 & TCP & Gestión administrativa Linux \\
\texttt{Allow-RDP} & 130 & Allow & 3389 & TCP & Acceso administrativo remoto Windows \\
\bottomrule
\end{tabularx}
\caption{Matriz de Control de Acceso y Filtrado en NSG.}
\end{table}

\begin{azureterminal}[az network nsg rule create (HTTP, HTTPS, SSH, RDP)]
$ az network nsg create \
    --resource-group "rg-bootcamp-part03-networking" \
    --name "nsg-bootcamp-sec" \
    --location "eastus"

$ az network nsg rule create -g rg-bootcamp-part03-networking --nsg-name nsg-bootcamp-sec \
    --name "Allow-HTTP" --priority 100 --direction Inbound --access Allow --protocol Tcp --destination-port-ranges 80

$ az network nsg rule create -g rg-bootcamp-part03-networking --nsg-name nsg-bootcamp-sec \
    --name "Allow-HTTPS" --priority 110 --direction Inbound --access Allow --protocol Tcp --destination-port-ranges 443

$ az network nsg rule create -g rg-bootcamp-part03-networking --nsg-name nsg-bootcamp-sec \
    --name "Allow-SSH" --priority 120 --direction Inbound --access Allow --protocol Tcp --destination-port-ranges 22

$ az network nsg rule create -g rg-bootcamp-part03-networking --nsg-name nsg-bootcamp-sec \
    --name "Allow-RDP" --priority 130 --direction Inbound --access Allow --protocol Tcp --destination-port-ranges 3389

ProvisioningState: Succeeded (Todas las reglas activas)
\end{azureterminal}

\subsection{3.3 Asignación de Interfaces de Red e IPs Públicas Estáticas}
Se crearon las interfaces de red dedicadas \texttt{nic-vm01} y \texttt{nic-vm02}, asociando una dirección IP pública estática (\texttt{pip-vm01}) bajo el estándar de resiliencia zonal \texttt{Standard SKU}, dejando la infraestructura lista para la vinculación segura de cargas de trabajo.

\begin{azureterminal}[az network public-ip create \& az network nic create]
$ az network public-ip create \
    --resource-group "rg-bootcamp-part03-networking" \
    --name "pip-vm01" \
    --sku Standard --allocation-method Static

$ az network nic create \
    --resource-group "rg-bootcamp-part03-networking" \
    --name "nic-vm01" \
    --vnet-name "VNet-01" \
    --subnet "Subnet-Frontend" \
    --network-security-group "nsg-bootcamp-sec" \
    --public-ip-address "pip-vm01"

$ az network nic create \
    --resource-group "rg-bootcamp-part03-networking" \
    --name "nic-vm02" \
    --vnet-name "VNet-01" \
    --subnet "Subnet-Backend" \
    --network-security-group "nsg-bootcamp-sec"

{
  "name": "nic-vm01",
  "provisioningState": "Succeeded",
  "ipConfigurations": [{ "privateIpAddress": "10.0.1.4" }]
}
\end{azureterminal}
